\chapter{二次型：合同、惯性与正定}\label{ch:quadratic}
\chaptermap{矩阵与合同}{化标准形三法}{正定四判据}{惯性与极值}
\begin{examtip}
\textbf{考研核心：}二次型矩阵、配方法、正交变换、标准形、规范形、惯性定理与正定判别。
\textbf{竞赛视角：}用正负子空间理解惯性，用Schur补处理分块正定，用谱界与配方解决不等式。合同保正负指数，通常不保特征值。
\end{examtip}
\section{二次型矩阵与惯性定理}
\subsection{标准形与规范形的区别}
\begin{definition}{二次型与合同}
实二次型是实变量的二次齐次多项式，零多项式也包含在内。它唯一表示为$f(\vx)=\vx\T A\vx$，其中对称矩阵$A=A\T$由平方项系数$a_{ii}$及交叉项系数$2a_{ij}$确定。作可逆代换$\vx=P\vy$后，矩阵变为$P\T AP$，称合同变换。对角形式$\sum d_i y_i^2$称标准形；系数只含$1,-1,0$时称规范形。
\end{definition}
\begin{theorem}{惯性定理与合同分类}
任何实对称矩阵均合同于$\diag(I_p,-I_q,0)$，且$p,q$与零项数不随可逆代换改变，称正、负惯性指数；秩$r=p+q$。两个同阶实对称矩阵合同，当且仅当其正负惯性指数分别相同。
\end{theorem}
\begin{proof}
\label{der:6-1}
实对称谱定理给$Q\T AQ=\diag(\lambda_i)$。对每个非零根作缩放$y_i=\frac{z_i}{\sqrt{|\lambda_i|}}$，零根方向保持不变，得到系数为$1,-1,0$的规范形；正交矩阵与这些非零缩放的乘积可逆，故为合同。

设一种规范形正系数个数为$p$，其正坐标子空间经可逆代换送回原向量空间，记为$U$；可逆代换保维数，故$\dim U=p$，且$f$在$U$上除零外严格正。若另一种规范形正系数为$p'$，将其非正坐标子空间也送回同一个原空间，记为$W$，则$\dim W=n-p'$，且$f$在$W$上不大于零。若$p>p'$，则$\dim U+\dim W>n$，由维数公式必有非零公共向量；在该向量上$f$既严格正又非正，矛盾。因此$p\le p'$，交换两种表示得$p'\le p$，故$p=p'$。对$-f$作相同论证得到负指数$q$也唯一。

可逆左右乘保持秩，规范形只有$p+q$个非零对角元，所以$\rank A=p+q$。若两实对称矩阵合同，惯性由上面论证相同；反之，它们分别合同于同一规范形$J$，设$P\T AP=J=Q\T BQ$，则$B=(PQ\inv)\T A(PQ\inv)$，因此合同。这里也验证了合同关系的反向构造。

\end{proof}
\begin{method}{先写对称矩阵再选变换}
逐项写对角系数和交叉系数的一半。合同前后秩相同，且$\det(P\T AP)=(\det P)^2\det A$；迹只在正交合同等相似变换下作不变性核验，不能用于一般合同。只求惯性时不必求全部特征值；若要求正交变换，则必须保留单位正交特征向量。易错：一般可逆配方与正交配方所得对角系数可不同。
\end{method}
\worked{交叉项的一半}{基础例题}{\star}{6-1}
{写出$f=2x^2+4xy+3y^2-2yz$的三阶对称矩阵。}
{注意没有$z^2$并不表示没有第三行列。}
{按变量次序$(x,y,z)$，平方项给$a_{11}=2,a_{22}=3,a_{33}=0$。交叉项$4xy$对应$2a_{12}=4$，所以$a_{12}=a_{21}=2$；$-2yz$对应$a_{23}=a_{32}=-1$；没有$xz$项，对应两个位置均0。因此
\[A=\begin{pmatrix}2&2&0\\2&3&-1\\0&-1&0\end{pmatrix}.\]
乘开$(x,y,z)A(x,y,z)\T$，交叉项为$2xy+2yx-yz-zy=4xy-2yz$，核验与题目一致。}
{交叉项的一半是二次型与矩阵之间最常见的转换失误。}
{将$4$直接填入$a_{12},a_{21}$会把交叉项翻倍。}
\worked{秩不等于正指数}{考研风格}{\star\star}{6-2}
{求$f=x^2-2y^2$的秩、惯性指数与规范形。}
{矩阵已对角，先读符号再缩放。}
{\textbf{方法一：}矩阵$\diag(1,-2)$有两个非零对角元，秩2；正系数一个、负系数一个，惯性$(1,1)$。令$u=x,v=\sqrt2y$，逆代换$x=u,y=\frac{v}{\sqrt2}$可逆，得$f=u^2-v^2$。
\textbf{方法二：}两特征值为1与$-2$，正负各一个，由正交谱表达得到同一惯性。正定子空间只能取一维，如$x$轴；整个平面含负方向$y$轴，不可能在全二维上正定。}
{秩给非零方向总数，正定方向最大维数是$p$。}
{不能由秩2推出存在二维正定子空间，本题恰是反例。}
\worked{合同不保特征值}{考研风格}{\star\star}{6-3}
{比较$A=\diag(1,2)$与$B=\diag(1,8)$的合同与相似关系。}
{合同允许非单位伸缩。}
{取可逆$P=\diag(1,2)$，逐项相乘得$P\T AP=\diag(1\cdot1^2,2\cdot2^2)=\diag(1,8)=B$，因此合同。相似则不成立：相似保迹，而$\tr A=3,\tr B=9$；也可比较谱$\{1,2\}\ne\{1,8\}$。两矩阵均正定，惯性都为$(2,0)$，与合同分类相符；合同的缩放会改变特征值，所以不能从合同反推相似。}
{合同关心符号与退化，相似关心线性变换本身。}
{正交合同同时也是相似，一般合同则不是。}
\worked{无平方项的二次型}{竞赛风格}{\star\star}{6-4}
{化$f=2xy$为标准形并求惯性。}
{用和差产生平方项。}
{\textbf{方法一：}令$u=\frac{x+y}{\sqrt2},v=\frac{x-y}{\sqrt2}$，逆代换$x=\frac{u+v}{\sqrt2},y=\frac{u-v}{\sqrt2}$。代入$2xy=(u+v)(u-v)=u^2-v^2$，规范形已得，惯性$(1,1)$。
\textbf{方法二：}矩阵$A=\begin{pmatrix}0&1\\1&0\end{pmatrix}$，$\chi_A=t^2-1$。和差向量分别给根1、$-1$，单位化拼$Q=\frac1{\sqrt2}\begin{pmatrix}1&1\\1&-1\end{pmatrix}$，$Q\T Q=I,Q\T AQ=\diag(1,-1)$，与第一法相同。秩为2，说明零对角不代表零秩。}
{缺平方项时先制造平方项，不能直接除零配方。}
{对角元全零不代表二次型为零，更不代表秩为零。}
\section{配方、正交与合同消元}
\subsection{同题多法及其输出差异}
\begin{theorem}{配方一步与合同消元}
若$a_{11}\ne0$，将$A=\begin{pmatrix}a&\vb\T\\\vb&C\end{pmatrix}$，则
\begin{equation}f=a\left(x_1+a^{-1}\vb\T\vz\right)^2+\vz\T(C-a^{-1}\vb\vb\T)\vz.\label{eq:complete}\end{equation}
这是可逆三角代换的合同。矩阵上对一列作操作时，再对对应行作同样操作，才形成$P\T AP$。若全部对角元为零但矩阵非零，先对非零$a_{ij}$对应变量作和差代换。
\end{theorem}
\begin{proof}
\label{der:6-2}
将$\vx=(x_1,\vz\T)\T$代入块矩阵，原式为
\[f=a x_1^2+2x_1\vb\T\vz+\vz\T C\vz.\]
因$a\ne0$，展开
\[a(x_1+a^{-1}\vb\T\vz)^2=a x_1^2+2x_1\vb\T\vz+a^{-1}\vz\T\vb\vb\T\vz.\]
从原式减去多出的末项，便得到式\eqref{eq:complete}。令$x_1=y_1-a^{-1}\vb\T\vz$，则变换矩阵
\[P=\begin{pmatrix}1&-a^{-1}\vb\T\\0&I\end{pmatrix},\quad\det P=1,
\quad P\T AP=\begin{pmatrix}a&0\\0&C-a^{-1}\vb\vb\T\end{pmatrix}.\]
这说明配方、成对行列消元与合同是同一个过程。若列操作由右乘初等矩阵$E$实现，对应的同行操作应左乘$E\T$，总变化为$E\T AE$。

若全部对角元为零但$a_{ij}\ne0$，仅在$i,j$两变量上令$x_i=u+v,x_j=u-v$，其交叉项变为$2a_{ij}(u^2-v^2)$，产生非零平方项；这个二阶代换行列式为$-2\ne0$。其余变量不变后即可继续上述配方，避免除以零对角元。

\end{proof}
\begin{method}{三法的目标差异}
配方得到任意标准形，计算短；正交法给特征值作系数并保长度；合同消元适合追踪变换矩阵。若要$P$，从$I$同步累积列操作。易错：只对行消元保秩，却一般不能保惯性。
\end{method}
\worked{同一二次型的三种化法}{考研风格}{\star\star}{6-5}
{将$f=2x^2+2xy+2y^2$化标准形，并比较变换。}
{既可配方，又可使用对称矩阵谱。}
{\textbf{方法一（配方）：}$2x^2+2xy=2[(x+\frac{y}{2})^2-\frac{y^2}{4}]$，所以$f=2(x+\frac{y}{2})^2+(\frac{3}{2})y^2$。设$u=x+\frac{y}{2},v=y$，逆代换$x=u-\frac{v}{2},y=v$，可逆，标准形为$2u^2+(\frac{3}{2})v^2$。
\textbf{方法二（正交）：}$A$的谱为3、1，取$Q=\frac1{\sqrt2}\begin{pmatrix}1&1\\1&-1\end{pmatrix}$。令$(x,y)\T=Q(u,v)\T$，直接代入可得$x^2+y^2=u^2+v^2,2xy=u^2-v^2$，故$f=3u^2+v^2$。
\textbf{方法三（合同消元）：}先$c_2\leftarrow c_2-\frac{c_1}{2}$，矩阵变为$\begin{pmatrix}2&0\\1&\frac{3}{2}\end{pmatrix}$；再$r_2\leftarrow r_2-\frac{r_1}{2}$，得到$\diag(2,\frac{3}{2})$。列操作矩阵$P=\begin{pmatrix}1&-\frac{1}{2}\\0&1\end{pmatrix}$，行操作为$P\T$，故确为$P\T AP$。两种标准形尺度不同，但正指数均2。}
{配方与合同消元是同一代换的代数和矩阵版本；正交法输出系数是特征值。}
{不同标准形不矛盾，正负指数相同即可；“标准形”不要求系数只能为$\pm1$。}
\worked{三元连续配方}{考研风格}{\star\star}{6-6}
{化$f=x^2+2xy+2y^2+2yz+2z^2$并给出可逆代换。}
{连续出现相邻交叉项。}
{先配$x$：$x^2+2xy+2y^2=(x+y)^2+y^2$。再配$y,z$：$y^2+2yz+2z^2=(y+z)^2+z^2$。所以$f=(x+y)^2+(y+z)^2+z^2$。令$u=x+y,v=y+z,w=z$，反解$z=w,y=v-w,x=u-v+w$，变换矩阵$P=\begin{pmatrix}1&-1&1\\0&1&-1\\0&0&1\end{pmatrix}$行列式1。因此标准形$u^2+v^2+w^2$，惯性$(3,0)$、秩3。每个新变量独立，未因配方丢失自由度。}
{写出逆代换排除了变量替换不独立的风险。}
{若题目要求正交变换，本代换不满足，需另求谱。}
\worked{全无平方项的三元例子}{竞赛风格}{\star\star\star}{6-7}
{化$f=2xy+2xz$为规范形。}
{把$y+z$看作一个方向，并补足第三坐标。}
{\textbf{方法一：}原式$2x(y+z)$。令$u=\frac{x+y+z}{\sqrt2},v=\frac{x-y-z}{\sqrt2},w=y-z$。则$u^2-v^2=2x(y+z)$；变换系数矩阵行列式2，非零，第三变量补足独立坐标。因此规范形$u^2-v^2$，惯性$(1,1)$、零项1、秩2。
\textbf{方法二：}$A=\begin{pmatrix}0&1&1\\1&0&0\\1&0&0\end{pmatrix}$，$\chi_A(t)=t(t^2-2)$，根$\sqrt2,-\sqrt2,0$。可取单位特征向量$q_+=\frac{(\sqrt2,1,1)\T}{2},q_-=\frac{(-\sqrt2,1,1)\T}{2},q_0=\frac{(0,1,-1)\T}{\sqrt2}$。逐点积为零、长度为1，故正交标准形$\sqrt2u^2-\sqrt2v^2$；再将两非零变量乘$2^{\frac{1}{4}}$即得同一规范形。}
{直接和差代换短，谱法同时给正交标准形。}
{$w$虽不出现在最终式中，仍须补入使整体代换可逆。}
\worked{完全平方与半正定}{考研风格}{\star\star}{6-8}
{研究$f=(x+y+z)^2$的秩、惯性与正交标准形。}
{矩阵是$\bm1\bm1\T$。}
{\textbf{方法一：}矩阵$A=\bm1\bm1\T$。$A\bm1=3\bm1$，和零平面上$Ax=0$。取$q_1=\frac{(1,1,1)\T}{\sqrt3},q_2=\frac{(1,-1,0)\T}{\sqrt2},q_3=\frac{(1,1,-2)\T}{\sqrt6}$。后两列分量和为零，故$Aq_2=Aq_3=0$；三列长度均1，点积的分子分别为$1-1+0=0,1+1-2=0,1-1+0=0$，所以$Q=(q_1,q_2,q_3)$正交。于是$AQ=Q\diag(3,0,0)$，左乘$Q\T$得$Q\T AQ=\diag(3,0,0)$，正交标准形$3u^2$，秩1、正指数1、负指数0。
\textbf{方法二：}令$u=x+y+z,v=y,w=z$，逆代换$x=u-v-w,y=v,z=w$可逆，原式直接变为$u^2$。它对所有向量非负，但例如$(1,-1,0)$非零且取值0，故半正定而非正定。}
{同一二次型的正交标准形系数3和规范形系数1对应不同缩放。}
{非零二次型也可能半正定而不正定，因有非零零值方向。}
\begin{insight}{正交变换是在更换单位正交坐标轴}
二次型$2x^2+2xy+2y^2$的等值线为倾斜椭圆。令$u=\frac{x+y}{\sqrt2},v=\frac{x-y}{\sqrt2}$，只改变观察坐标，椭圆的几何形状与距离都不变，交叉项却消失，变成$3u^2+v^2=1$。特征值3对应短半轴$\frac{1}{\sqrt3}$，特征值1对应长半轴1。本例坐标变换矩阵行列式为$-1$，含方向反转；一般正交变换包括旋转与反射，不都属于纯旋转。
\end{insight}
\begin{figure}[H]\centering
\begin{tikzpicture}[scale=1.6]
\begin{scope}
\draw[laaxis](-1.35,0)--(1.4,0)node[right]{$x$};\draw[laaxis](0,-1.35)--(0,1.4)node[above]{$y$};
\draw[Teal,thick,rotate=45](0,0)ellipse({1/sqrt(3)} and 1);
\draw[dashed,Gold](-.85,-.85)--(.85,.85);\draw[dashed,Gold](-.85,.85)--(.85,-.85);
\node[lanote]at(0,-1.75){$2x^2+2xy+2y^2=1$};
\end{scope}
\draw[-{Stealth},Ink,thick](1.65,.15)--(2.65,.15)node[midway,above]{$Q\T$};
\begin{scope}[xshift=4.4cm]
\draw[laaxis](-1.35,0)--(1.4,0)node[right]{$u$};\draw[laaxis](0,-1.35)--(0,1.4)node[above]{$v$};
\draw[Teal,thick](0,0)ellipse({1/sqrt(3)} and 1);
\draw[Gold,thick](0,0)--({1/sqrt(3)},0)node[below right]{$\frac{1}{\sqrt3}$};
\draw[Gold,thick](0,0)--(0,1)node[above left]{$1$};
\node[lanote]at(0,-1.75){$3u^2+v^2=1$};
\end{scope}
\end{tikzpicture}
\caption{沿主轴观察，交叉项消失；正交变换保留半轴长度。}\label{fig:quadratic-ellipse}
\end{figure}
\begin{clevertrick}{二阶正定有时只需“迹与行列式”}
对实对称二阶矩阵，谱定理保证两个特征值$\lambda_1,\lambda_2$都实。若$\det A=\lambda_1\lambda_2>0$，两根同号且均非零；再由$\tr A=\lambda_1+\lambda_2>0$排除都负，只能全正，故$A\succ0$。反过来正定使两根都正，其和、积自然正，因此
\[A\succ0\quad\Longleftrightarrow\quad\det A>0\ \text{且}\ \tr A>0.\]
例如$A=\begin{pmatrix}t&1\\1&t\end{pmatrix}$，行列式$t^2-1$、迹$2t$。前一不等式给$t<-1$或$t>1$，后一给$t>0$；取交集得$t>1$，端点$t=1$因出现零根必须剔除。
\textbf{易错提示：}这条捷径依赖二阶和实对称；三阶矩阵$\diag(-1,-1,3)$迹、行列式都正，却不正定。
\end{clevertrick}


\section{正定与半正定的判据}
\subsection{四种正定判定必须满足条件}
\begin{definition}{正定、负定与不定}
实对称$A$称正定若对每个$\vx\ne0$有$\vx\T A\vx>0$；称半正定若对每个$\vx$有$\vx\T A\vx\ge0$。负定、半负定分别将符号改为$<$、$\le$。本册半正定包含正定，半负定包含负定。不定是指存在$u,v$使$u\T Au>0$且$v\T Av<0$。一个非零向量给零值，本身不能推出不定。

取同一非空指标集$I\subseteq\{1,\ldots,n\}$同时选行和列，$\det A[I,I]$称主子式；特别取$I=\{1,\ldots,k\}$所得$\Delta_k=\det A[1\!:\!k,1\!:\!k]$称$k$阶顺序主子式。只取相同行号与列号才是主子式，顺序主子式只是其中的一部分。
\end{definition}
\begin{theorem}{正定四判据与半正定判据}
实对称矩阵下列等价：定义上的正定；所有特征值正；各阶顺序主子式$\Delta_k>0$；合同规范形为$I_n$。又等价于存在可逆$R$使$A=R\T R$。负定等价于$(-1)^k\Delta_k>0$。半正定等价于全部特征值非负，也等价于\textbf{所有主子式}非负，不能只检查顺序主子式。
\end{theorem}
\begin{proof}
\label{der:6-3}
\textbf{定义、谱与规范形。}令$x=Qy$，$Q$正交，则$f=\sum_i\lambda_i y_i^2$。若全$\lambda_i>0$，非零$y$至少一项正，故正定；若某$\lambda_i\le0$，取对应$y=e_i$便否定正定。故正定等价于全正谱，也等价于惯性$(n,0)$。令$R=\diag(\sqrt{\lambda_i})Q\T$，则$A=R\T R$且$R$可逆；反之$R$可逆时$x\T R\T Rx=\norm{Rx}^2>0$。

\textbf{顺序主子式。}正定矩阵限制到前$k$坐标仍正定，故该子矩阵的特征值全正，行列式$\Delta_k>0$。反向以维数归纳：$n=1$明显。若全部$\Delta_k>0$，则$a_{11}=\Delta_1>0$；按推导\ref{der:6-2}合同到$\diag(a_{11},S)$。$S$的前$k$阶顺序主子式为$\frac{\Delta_{k+1}}{a_{11}}>0$，这是对相应前$k+1$阶块应用Schur行列式公式所得。归纳知$S$正定，故整式正定。连续主元相应为$d_k=\frac{\Delta_k}{\Delta_{k-1}}$，$\Delta_0=1$。负定就是$-A$正定，$-A$的$k$阶顺序主子式为$(-1)^k\Delta_k$，得到负定判据。

\textbf{半正定与所有主子式。}若半正定，任一主子矩阵也半正定，谱非负，其行列式非负。反向，设全部主子式非负。任一$k$阶主子矩阵$B$满足
\[\det(B+\varepsilon I_k)=\sum_{J\subseteq\{1,\ldots,k\}}\varepsilon^{k-|J|}\det B[J,J].\]
此式由列多线性得到：选择$\varepsilon e_j$的列占据同号的对应行，未选位置留下同指标主子式，空主子式取1。各系数非负且$\varepsilon^k>0$，所以对$\varepsilon>0$全部顺序主子式严格正，$A+\varepsilon I$正定。固定$x$有$x\T Ax+\varepsilon\norm x^2>0$，令$\varepsilon\downarrow0$得到$x\T Ax\ge0$。这证明充分性，并说明为何只检查顺序主子式非负不够。

\end{proof}
\begin{method}{正定四条路怎么选}
低阶参数题用顺序主子式；结构矩阵用谱；能配方时用正平方和；证明题用定义与可逆合同。半正定优先用谱或平方和；只检验几个对角元或行列式不够。
\end{method}
\subsection{应试反例：非对称矩阵不能直接查顺序主子式}\label{r2:positive}
本章的 Sylvester 正定判据明确以\textbf{实对称矩阵}为对象。对任意实方阵 $A$，令 $S=(A+A^{\mathsf T})/2$、$K=(A-A^{\mathsf T})/2$，则 $K^{\mathsf T}=-K$，且
\[
 x^{\mathsf T}Kx=0,\qquad x^{\mathsf T}Ax=x^{\mathsf T}Sx.
\]
故对非对称系数阵，二次式的正负应检查 $S$。按本书约定，“半正定矩阵”本身也要求对称；若不对称，应说“它定义的二次式非负”或“它的对称部分半正定”。

\begin{example}{顺序主子式全正却不产生正定二次式}{考研辨析}{r2-positive}
取 $A=\begin{pmatrix}1&2\\0&1\end{pmatrix}$，两阶顺序主子式均为一，两个特征值也都是一。但
\[
 S=\frac{A+A^{\mathsf T}}2=\begin{pmatrix}1&1\\1&1\end{pmatrix},\qquad
 x^{\mathsf T}Ax=(x_1+x_2)^2.
\]
非零向量 $(1,-1)^{\mathsf T}$ 使二次式为零，因此它不严格为正。$S$ 的特征值为 $0,2$，是半正定的实对称矩阵，而 $A$ 本身不在本书正定/半正定矩阵定义的对象类中。
\end{example}

\begin{warning}
先查实对称，再选谱、主子式、配方或合同法。对于实对称矩阵，严格正定看全部顺序主子式正；半正定看所有主子式非负。不能同时省掉“对称”和“所有”两个条件。
\end{warning}

\paragraph{自测与答案}
$B=\begin{pmatrix}1&3\\0&1\end{pmatrix}$ 的顺序主子式也都为一，但 $x=(1,-1)^{\mathsf T}$ 给 $x^{\mathsf T}Bx=-1$，而 $x=(1,0)^{\mathsf T}$ 给一，所以其二次式不定。正特征值对一般非对称矩阵也不是二次式正定的充分条件。


\worked{参数正定的三种判定}{考研风格}{\star\star}{6-9}
{求$A_t=\begin{pmatrix}t&1\\1&t\end{pmatrix}$正定的参数范围。}
{对称二阶题可三路对照。}
{\textbf{方法一：}顺序主子式$\Delta_1=t,\Delta_2=t^2-1$，正定要求$t>0$且$t<-1$或$t>1$，交集为$t>1$。
\textbf{方法二：}和差特征值为$t+1,t-1$，两者全正恰当于$t>1$。
\textbf{方法三：}先用向量$(1,0)$看出正定必$t>0$，随后合法配方
$f=t(x+\frac{y}{t})^2+(t-\frac{1}{t})y^2$。可逆三角代换下两系数须正，即$t>0,t^2>1$，仍得$t>1$。边界$t=1$有零方向，不能包含进严格正定范围。}
{谱法最短；主子式最通用；配方法需先分$t=0$的边界。}
{$t=1$为半正定，$t=-1$为半负定，均非正定。}
\worked{顺序主子式不足以判半正定}{基础例题}{\star\star}{6-10}
{举例说明顺序主子式全非负不保证半正定。}
{让首个主元为零，同时隐藏一个负方向。}
{取$A=\diag(0,-1)$。第一阶顺序主子式$a_{11}=0$，第二阶$\det A=0$，二者均非负；但$x=e_2$给$x\T Ax=-1<0$，所以不是半正定。若检查所有主子式，除了$a_{11}$还必须检查$a_{22}=-1$，立即发现负值。这一反例明确否定“顺序主子式非负足够”的错误替换。}
{严格正与非负的判据不能简单换符号。}
{半正定需要所有主子式，不能把“顺序”二字漏掉后仍只算前几个。}
\worked{非对称矩阵的二次式}{竞赛风格}{\star\star}{6-11}
{研究$A=\begin{pmatrix}1&10\\-10&1\end{pmatrix}$的$\vx\T A\vx$，并说明与正定对称矩阵的关系。}
{反对称部分在二次式中消失。}
{\textbf{方法一：}逐项展开$x\T Ax=x^2+10xy-10yx+y^2=x^2+y^2$，对非零$(x,y)$严格正。
\textbf{方法二：}分解$A=S+K$，$S=\frac{A+A\T}{2}=I$，$K=\frac{A-A\T}{2}$反对称。标量$s=x\T Kx$满足$s=s\T=x\T K\T x=-s$，故$s=0$，二次式只剩$x\T Sx$。按本册约定正定矩阵本身要求对称，因此准确说法是$A$的对称部分正定；不能把非对称$A$直接代入实对称判据。}
{二次式只反映对称部分，不能从中恢复全部矩阵。}
{非对称$A$的特征值为$1\pm10i$，不能套实对称谱判据。}
\worked{Schur补判正定}{竞赛风格}{\star\star\star}{6-12}
{求$A=\begin{pmatrix}2&1&0\\1&2&1\\0&1&t\end{pmatrix}$正定的充要条件。}
{前二阶块已经正定。}
{\textbf{方法一：}$\Delta_1=2$，$\Delta_2=4-1=3$；沿末行或三对角递推得$\Delta_3=t\cdot3-1\cdot2=3t-2$。前两者已正，故正定恰当于$t>\frac{2}{3}$。
\textbf{方法二：}以前二阶块$B=\begin{pmatrix}2&1\\1&2\end{pmatrix}\succ0$，耦合列$c=(0,1)\T$。$B\inv=\frac13\begin{pmatrix}2&-1\\-1&2\end{pmatrix}$，所以$c\T B\inv c=\frac{2}{3}$。配方将二次型化成$u\T Bu+(t-\frac{2}{3})z^2$，可逆代换保持正定，因此仍要求$t>\frac{2}{3}$。这也直接说明边界会出现一个零方向。}
{主子式计算短；Schur补适合更大的分块矩阵。}
{一般Schur正定判据要求主块实对称且正定，只可逆还不够。}
\section{惯性、极值与平方根}
\subsection{把结构用于证明和不等式}
\begin{corollary}{正负子空间与平方根}
二次型正定子空间的最大维数为$p$，负定子空间的最大维数为$q$。正定$A=Q\Lambda Q\T$有唯一正定对称平方根$S=Q\diag(\sqrt{\lambda_i})Q\T$。实对称矩阵一般可写为带符号外积和；仅半正定时能全部使用正号。
\end{corollary}
\begin{proof}
\label{der:6-4}
规范形中正坐标空间给出$p$维正定子空间。任意维数大于$p$的子空间，与维数$n-p$的非正坐标空间维数和大于$n$，有非零交向量，在该向量上二次型不可能严格正。因此最大正定子空间维数恰为$p$；对$-A$同样得到最大负定维数为$q$。

若$A=Q\Lambda Q\T\succ0$，定义$S=Q\diag(\sqrt{\lambda_i})Q\T$。它实对称，谱全正，故正定；相乘时$Q\T Q=I$，所以$S^2=Q\Lambda Q\T=A$。若另有正定对称$T$满足$T^2=A$，则$TA=T^3=AT$，故$T$保留$A$的每个特征空间。对$A$的$\lambda$特征空间，$T$的限制仍正定对称且平方为$\lambda I$；正交对角化该限制，所有特征值$\mu$都满足$\mu^2=\lambda,\mu>0$，只能为$\sqrt\lambda$，所以限制恰为$\sqrt\lambda I$。各特征空间直和为全空间，故$T=S$，证明唯一。

最后，谱分解写为
\[A=\sum_{\lambda_i\ne0}\sgn(\lambda_i)\bigl(\sqrt{|\lambda_i|}q_i\bigr)\bigl(\sqrt{|\lambda_i|}q_i\bigr)\T.\]
若全用正号，任意$x$代入后都是平方和，必半正定；反向半正定时全部非零谱值正，上式全为正号。由此外积表示的准确条件得到证明。

\end{proof}
\begin{method}{极值先检查约束的几何}
单位球上的二次型用最小、最大特征值；线性约束下可先消元配方；正定权约束可用平方根变成单位球。要求规范形只需符号，要求数值极值则必须保留尺度。
\end{method}
\worked{含参数的惯性分类}{考研风格}{\star\star}{6-13}
{求$f=x^2+2xy+y^2+tz^2$的惯性和秩。}
{先把前两项合成平方。}
{配方$f=(x+y)^2+tz^2$。令$u=x+y,v=x-y,w=z$，前二阶变换行列式$-2$，所以整体可逆；$f=u^2+tw^2$，$v$为零方向。$t>0$时两正项，$(p,q,r)=(2,0,2)$；$t=0$时只剩一正项，$(1,0,1)$；$t<0$时一正一负，$(1,1,2)$。也可将矩阵在和、差、$z$方向正交化，其谱为$2,0,t$，逐符号计数与配方完全一致。}
{配方避开求三阶特征多项式。}
{三个变量不等于秩3，零项方向必须计入。}
\worked{加权约束下的最小值}{竞赛风格}{\star\star\star}{6-14}
{求$x+y=1$时$2x^2+y^2$的最小值。}
{配方或加权Cauchy都适合。}
{\textbf{方法一：}代入$y=1-x$，得$2x^2+(1-x)^2=3x^2-2x+1=3(x-\frac{1}{3})^2+\frac{2}{3}$，最小值$\frac{2}{3}$在$x=\frac{1}{3},y=\frac{2}{3}$取得。
\textbf{方法二：}写$x+y=(\sqrt2x)(\frac{1}{\sqrt2})+y\cdot1$，Cauchy给
\[1\le(2x^2+y^2)(\dfrac{1}{2}+1)=\dfrac32(2x^2+y^2).\]
因此目标至少$\frac{2}{3}$。等号要求$(\sqrt2x,y)$与$(\frac{1}{\sqrt2},1)$成比例，即$y=2x$；再结合$x+y=1$得到同一点，证明下界确达到且最优点唯一。}
{权重改变了最优方向，不能直接令$x=y$。}
{Cauchy等号条件需与原约束一起求解。}
\worked{正定平方根的显式计算}{竞赛风格}{\star\star}{6-15}
{求$A=\begin{pmatrix}2&1\\1&2\end{pmatrix}$的正定对称平方根。}
{使用已知和差特征方向。}
{\textbf{方法一：}以单位和差向量拼$Q$，$A=Q\diag(3,1)Q\T$，所以
\[S=Q\diag(\sqrt3,1)Q\T=\dfrac12\begin{pmatrix}\sqrt3+1&\sqrt3-1\\\sqrt3-1&\sqrt3+1\end{pmatrix}.\]
谱$\sqrt3,1$全正且$S$对称，相乘$S^2=A$，正定平方根由唯一性定理确定。
\textbf{方法二：}由对称结构设候选$S=\begin{pmatrix}a&b\\b&a\end{pmatrix}$。平方要求$a^2+b^2=2,2ab=1$，故$(a+b)^2=3,(a-b)^2=1$。正定要求特征值$a+b,a-b$均正，取$a+b=\sqrt3,a-b=1$，得$a=\frac{\sqrt3+1}{2},b=\frac{\sqrt3-1}{2}$。此构造确满足条件，再由唯一性排除其他形状的正定平方根。}
{谱法自然保证正定；待定系数法须筛选符号。}
{矩阵平方根不止一个；唯一性要求正定且对称。}
\worked{矩形变换后的正定条件}{竞赛风格}{\star\star\star}{6-16}
{设$A\succ0$、$B\in\R^{m\times n}$，其中$A$为$m$阶。证明$B\T AB\succ0$当且仅当$B$满列秩。}
{考察$B\vx$是否可能对非零$\vx$变成零。}
{\textbf{方法一：}对任意$x\ne0$，$x\T B\T ABx=(Bx)\T A(Bx)$。若$B$满列秩，$Bx\ne0$，正定$A$使其严格正；若$B$不满列秩，存在非零$x$使$Bx=0$，该二次式为0，不能正定，证明充要。
\textbf{方法二：}用正定平方根$A=S\T S$，且$S$可逆，得$B\T AB=(SB)\T(SB)$。Gram正定当且仅当$SB$满列秩；可逆左乘不改秩，$\rank SB=\rank B$，故仍等价于$\rank B=n$。无论是否满列秩，该式至少半正定。}
{定义法最短，平方根法联系矩阵分解。}
{无论$B$是否满列秩，该矩阵都半正定；严格正需要额外条件。}
\begin{insight}{惯性决定等值线的基本类型}
二维正定型的正等值线是椭圆；一个正方向和一个负方向给双曲线；一个正方向加一个零方向，则图形沿零方向无限延伸。合同可以改变长度尺度，却不能把正方向变成负方向，也不能凭可逆代换消除零方向。
\end{insight}
\begin{figure}[H]\centering
\begin{tikzpicture}[scale=1.1]
\begin{scope}
\draw[laaxis](-1.4,0)--(1.5,0);\draw[laaxis](0,-1.3)--(0,1.4);
\draw[Teal,thick](0,0)ellipse(1 and {1/sqrt(2)});\node[lanote]at(0,-1.85){椭圆：$p=2,q=0$\\$x^2+2y^2=1$};
\end{scope}
\begin{scope}[xshift=4.3cm]
\draw[laaxis](-1.4,0)--(1.5,0);\draw[laaxis](0,-1.3)--(0,1.4);
\draw[Teal,thick,domain=-.9:.9,samples=60,variable=\t]plot({sqrt(1+\t*\t)},{\t});
\draw[Teal,thick,domain=-.9:.9,samples=60,variable=\t]plot({-sqrt(1+\t*\t)},{\t});
\node[lanote]at(0,-1.85){双曲线：$p=q=1$\\$x^2-y^2=1$};
\end{scope}
\begin{scope}[xshift=8.6cm]
\draw[laaxis](-1.4,0)--(1.5,0);\draw[laaxis](0,-1.3)--(0,1.4);
\draw[Teal,thick](-1,-1.1)--(-1,1.1);\draw[Teal,thick](1,-1.1)--(1,1.1);
\node[lanote]at(0,-1.85){平行线：$p=1,q=0$\\$x^2=1$，一个零方向};
\end{scope}
\end{tikzpicture}
\caption{正、负、零方向决定类型；惯性与图形开闭特征一致。}\label{fig:inertia-contours}
\end{figure}
\begin{clevertrick}{正定型的线性约束可用加权Cauchy}
对实对称$A\succ0$，取其对称正定平方根$S=A^{\frac{1}{2}}$。令$\vy=S\vx,w=S\inv\vc$，则目标$\vx\T A\vx=\norm\vy^2$，约束$\vc\T\vx=w\T\vy=b$，且$\norm w^2=\vc\T A\inv\vc>0$。Cauchy给
\[b^2=(w\T\vy)^2\le\norm w^2\norm\vy^2=(\vc\T A\inv\vc)(\vx\T A\vx),\]
除以正分母便得目标下界。等号要求$\vy=\alpha w$，代入约束得$\alpha\norm w^2=b$，所以唯一等号向量为
\[\vx_{\min}=S\inv\vy=\dfrac{b}{\vc\T A\inv\vc}A\inv\vc,\qquad
\min\vx\T A\vx=\dfrac{b^2}{\vc\T A\inv\vc}.\]
这里$\vc\ne0$，使$w\ne0$并保证分母严格正；$b=0$时最优向量为零，同样包含。对$2x^2+y^2$与$x+y=1$，$A=\diag(2,1),\vc=(1,1)\T$，故$A\inv\vc=(\frac{1}{2},1)\T$、$\vc\T A\inv\vc=\frac{3}{2}$。最优列向量为$\vx_* =\frac{2}{3}(\frac{1}{2},1)\T=(\frac{1}{3},\frac{2}{3})\T$，最小值$\frac{1}{\frac{3}{2}}=\frac{2}{3}$；代回$\frac{2}{9}+\frac{4}{9}=\frac{2}{3}$，核验成立。
\textbf{易错提示：}正定保证椭圆闭合与分母严格正；不定型不能无条件套用这一最小值公式。
\end{clevertrick}


\begin{warning}
正定不是“行列式正”，例如$-I_2$行列式正却负定。半正定不是“顺序主子式非负”。一般秩$r$的二次型不一定合同于$\diag(I_r,0)$；它还可能有负指数。可逆代换保证惯性不变，非可逆代换可能删掉方向。
\end{warning}
\section*{章末分层习题}\addcontentsline{toc}{section}{章末分层习题}
\exercise{6-1}{基础}{写出$x^2+6xy+4y^2$的矩阵。}
\exercise{6-2}{基础}{求$2x^2-3y^2$的规范形和惯性。}
\exercise{6-3}{基础}{判断$\diag(1,0,2)$的正定性、半正定性和秩。}
\exercise{6-4}{考研风格}{用配方和正交两法化$3x^2+2xy+3y^2$。}
\exercise{6-5}{考研风格}{求$\begin{pmatrix}t&2\\2&3\end{pmatrix}$正定的参数范围。}
\exercise{6-6}{考研风格}{判断$\diag(-1,-2,3)$的惯性与正定子空间的最大维数。}
\exercise{6-7}{竞赛拓展}{证明实对称半正定$A$可表示成$r=\rank A$个非零向量外积之和。}
\exercise{6-8}{竞赛拓展}{求$x^2+y^2=1$时$3x^2+2xy+3y^2$的最值。}
\exercise{6-9}{竞赛拓展}{设$A\succ0$，证明$\det(A+\vu\vu\T)>\det A$对非零$\vu$成立。}
\section*{章末方法速查}\addcontentsline{toc}{section}{章末方法速查}
\begin{quicktable}
配方&不要求正交&交叉项可快速消去&代换必须可逆\\
正交变换&实对称&要求保长度或谱&单位化与列次序\\
合同消元&实对称&需要变换矩阵&行列成对操作\\
顺序主子式&正定、负定&低阶含参数&半正定不能照搬\\
惯性分类&合同判断&只需符号&秩不等于正指数\\
Schur补&正定主块&分块与带状&主块先正定\\
Rayleigh商&单位球约束&二次式极值&规范形丢尺度\\
\end{quicktable}
